\chapter{CPU Microarchitecture}\label{ch:ll:cpu-microarchitecture}

In June 2012 a programmer asked why twelve lines of C++, adding up the elements of an array that are at least 128,
ran in 11.5 seconds on random data and 1.9 seconds once the same data had been sorted. Nothing in the source had
changed; only the order of the values, and with it how often the processor guessed correctly which way an
\texttt{if} would go. The answer became one of the most read in the history of the site. A trading system's hot path
is a few thousand instructions long, and the processor that runs them is not the sequential machine the source code
describes: it overlaps dozens of instructions, executes them out of order, guesses at branches, and changes its own
clock speed. This chapter describes that machine from the program's side, measures each effect on this book's laptop,
and builds the harness every later chapter times itself with.

\section{Pipelines and superscalar execution}

\begin{definition}[Instruction pipeline, instructions per cycle]\label{def:ll:cpu-microarchitecture:pipeline}
An \emph{instruction pipeline}\index{instruction pipeline} splits the execution of an instruction into stages (fetch,
decode, rename, schedule, execute, retire), each handled by its own hardware, so that many instructions are in
flight at once, each in a different stage. A superscalar pipeline handles several instructions per stage per cycle.
\emph{Instructions per cycle}\index{instructions per cycle} (IPC) is the number of instructions a core completes in a
cycle of its clock, averaged over a piece of code.
\end{definition}

A current performance core decodes and renames six instructions a cycle and can sustain six instructions a cycle
retired, when the instructions are independent. Code rarely comes close: an IPC between one and three is ordinary.
The difference between the width of the machine and the IPC achieved is where latency engineering lives: every
cycle in which fewer instructions retire than could is a cycle spent waiting, for a result, for memory, or for the
front of the pipeline to recover from a wrong guess.

\begin{omfigure}
\begin{tikzpicture}[font=\footnotesize,
  st/.style={draw=omDef,fill=omDef!8,rounded corners=2pt,minimum height=0.8cm,minimum width=1.55cm,align=center}]
\node[st] (f) at (0,0) {fetch\\ (predict)};
\node[st] (d) at (2.1,0) {decode};
\node[st] (r) at (4.2,0) {rename};
\node[st] (s) at (6.3,0) {schedule};
\node[st,minimum height=2.0cm] (x) at (8.5,0) {execute\\ (ALUs,\\ loads,\\ stores)};
\node[st] (t) at (10.7,0) {retire};
\foreach \a/\b in {f/d,d/r,r/s,s/x,x/t} \draw[-Stealth,thick] (\a) -- (\b);
\draw[draw=omThm,fill=omThm!8,rounded corners=2pt] (3.4,-1.9) rectangle (11.5,-1.3);
\node[font=\footnotesize] at (7.45,-1.6) {reorder buffer: program order, completed out of order};
\draw[-Stealth,omThm] (r.south) -- (4.2,-1.3);
\draw[-Stealth,omThm] (10.7,-1.3) -- (t.south);
\draw[-Stealth,thick,omProp] (x.north) to[out=120,in=60,looseness=0.6] node[pos=0.5,above,font=\footnotesize]{mispredicted branch: flush and refetch} (f.north);
\end{tikzpicture}

\omcaption{The pipeline of an out-of-order core, simplified. Instructions enter and leave in program order; in between,
the scheduler starts each one as soon as its inputs are ready, and the \omterm{def:ll:cpu-microarchitecture:ooo}{reorder buffer} keeps the program's order for
retirement. A mispredicted branch discards every instruction fetched after it.}\label{fig:ll:cpu-microarchitecture:pipeline}
\end{omfigure}

\section{Out-of-order execution and the reorder buffer}

\begin{definition}[Out-of-order execution, reorder buffer]\label{def:ll:cpu-microarchitecture:ooo}
In \emph{out-of-order execution}\index{out-of-order execution} the core starts each instruction as soon as its inputs
are available rather than in program order, and makes the results visible in program order. The
\emph{reorder buffer}\index{reorder buffer} holds the instructions between rename and retirement; its size bounds how
far ahead of a stalled instruction the core can find independent work.
\end{definition}

The \omterm{def:ll:cpu-microarchitecture:ooo}{reorder buffer} of the performance core family in this laptop holds 512 instructions. When a load misses every cache
and waits a hundred nanoseconds for memory (chapter~3), the core keeps executing the independent instructions behind
it until the buffer is full, then stalls. What it cannot do is shorten a dependency chain: each instruction that
needs the previous one's result waits for it, whatever the width of the machine.

To see these effects the book times small pieces of code in isolation.

\begin{definition}[Microbenchmark, compiler barrier]\label{def:ll:cpu-microarchitecture:ubench}
A \emph{microbenchmark}\index{microbenchmark} times a small piece of code in isolation, many times, to measure one
property of the machine or of the code. A \emph{compiler barrier}\index{compiler barrier} is a statement the compiler
must treat as reading or writing memory (or a given value) although it emits no instruction; it stops the optimiser
from deleting or moving the code being timed.
\end{definition}

A \omterm{def:ll:cpu-microarchitecture:ubench}{microbenchmark} measures what it measures: a loop over data that stays in the first-level cache, on a warm core,
after its branches have been learnt. That is the best case of the hot path, not its behaviour after a quiet period
(cold caches, a core in a low-power state), which is often what matters in trading (chapter~5).

\begin{proposition}[Chains needed to saturate the units]\label{prop:ll:cpu-microarchitecture:little}
If an operation has latency $\ell$ cycles and the core can start $w$ of them per cycle, a loop needs at least $\ell w$
independent dependency chains to start $w$ a cycle; with $k < \ell w$ chains it completes at most $k/\ell$ a cycle.
\end{proposition}

\begin{proof}
Each chain has at most one operation in flight, occupying $\ell$ cycles; $k$ chains have at most $k$ in flight. By
Little's law the completion rate is at most $k/\ell$, which reaches $w$ only when $k \ge \ell w$.
\end{proof}

\Cref{fig:ll:cpu-microarchitecture:chains} sums 4\,096 doubles, held in cache, with one to sixteen accumulators. A
floating-point addition on this core has a latency of two cycles when followed by another addition, and two can start
per cycle: \cref{prop:ll:cpu-microarchitecture:little} predicts four chains, and the measured cost per addition halves
from one accumulator to two, halves again to four, and stops improving there. The compiler does not do this itself:
without permission to reassociate floating-point additions (\texttt{-ffast-math}, chapter~8) it must keep the one
chain the source wrote.

\begin{omfigure}
\begin{tikzpicture}
\begin{axis}[width=9cm,height=5.2cm,ybar,bar width=14pt,ymin=0,enlarge y limits={upper,value=0.2},
  symbolic x coords={1,2,4,8,16},xtick=data,enlarge x limits=0.15,
  xlabel={independent accumulators},ylabel={ns per addition},tick label style={font=\small},label style={font=\small},
  grid=major,grid style={gray!25},nodes near coords,nodes near coords style={font=\footnotesize},
  point meta=y,every node near coord/.append style={/pgf/number format/fixed,/pgf/number format/precision=2},
  table/col sep=comma]
\addplot[fill=omDef!55,draw=omDef] table[x=accumulators,y=ns_per_add]{figdata/low-latency/02-cpu-microarchitecture/measured_chains.csv};
\end{axis}
\end{tikzpicture}

\omcaption{Cost per floating-point addition when summing 4\,096 doubles with 1 to 16 independent accumulators. Two
cycles of latency and two adders make four chains enough (\cref{prop:ll:cpu-microarchitecture:little}). Measured on a
laptop (Intel Core Ultra 7 155H) under WSL2, no isolated cores, pinned to one CPU, best of 30 runs. Data:
\texttt{bench\_cpu.py}.}\label{fig:ll:cpu-microarchitecture:chains}
\end{omfigure}

The core can also remove dependencies the program contains. A chain of ten dependent additions of a register runs at
one addition a cycle, as it should; the same chain adding the constant 1 ran at about five additions a cycle on this
laptop, because the renamer of this core family resolves small-immediate additions itself, up to six a cycle, without
sending them to an execution unit. A \omterm{def:ll:cpu-microarchitecture:ubench}{microbenchmark} that uses \texttt{add \$1} to count cycles measures the renamer;
the harness below measures the clock with the register form.

\section{Branch prediction}

\begin{definition}[Branch prediction, branch misprediction]\label{def:ll:cpu-microarchitecture:branch}
\emph{Branch prediction}\index{branch prediction} is the core's guess, made at fetch, of the direction and target of
each branch, so that the pipeline can keep fetching before the branch is executed. A \emph{branch
misprediction}\index{branch misprediction} is a wrong guess: the core discards every instruction fetched after the
branch and refetches from the right target, losing the depth of the front of the pipeline in cycles.
\end{definition}

\begin{proposition}[When branchless code wins]\label{prop:ll:cpu-microarchitecture:breakeven}
A loop body with a data-dependent branch costs $c_0 + mP$ cycles per iteration on average, where $m$ is the
misprediction rate and $P$ the misprediction penalty; a branchless rewrite that always does the work of both sides
costs $c_0 + c$. The rewrite is cheaper exactly when $m > c/P$.
\end{proposition}

\begin{proof}
Each iteration mispredicts with probability $m$ and then pays $P$; linearity of expectation gives $c_0 + mP$. Compare
with $c_0 + c$.
\end{proof}

\Cref{fig:ll:cpu-microarchitecture:branch} reruns the 2012 loop. On sorted data the predictor is right on all but one
of 32\,768 iterations and the branchy loop is the fastest; on shuffled data it is right half the time and the loop is
about thirteen times slower. The branchless loop, which masks each element with all ones or all zeros, costs the same
on both. Half the iterations mispredict on shuffled data, so the difference between the two branchy bars is half the
penalty per element: about \qty{5.2}{\nano\second}, or some 24 cycles at this core's measured clock, a little above the
``sometimes exceeding 20 clock cycles'' a published measurement gives for this core family. The branchless loop
costs about \qty{0.11}{\nano\second} more per element on sorted data, so by
\cref{prop:ll:cpu-microarchitecture:breakeven} it wins as soon as more than about two iterations in a hundred
mispredict.

\begin{omfigure}
\begin{tikzpicture}
\begin{axis}[width=9cm,height=5.2cm,ybar,bar width=16pt,ymin=0,
  symbolic x coords={sorted,shuffled},xtick=data,enlarge x limits=0.5,
  ylabel={ns per element},tick label style={font=\small},label style={font=\small},
  grid=major,grid style={gray!25},
  legend columns=2,legend style={font=\footnotesize,at={(0.5,-0.2)},anchor=north,draw=none,fill=none,
  /tikz/every even column/.append style={column sep=8pt}},table/col sep=comma]
\addplot[fill=omThm!45,draw=omThm] table[x=order,y=branchy]{figdata/low-latency/02-cpu-microarchitecture/measured_branch.csv};
\addplot[fill=omDef!45,draw=omDef] table[x=order,y=branchless]{figdata/low-latency/02-cpu-microarchitecture/measured_branch.csv};
\legend{branchy,branchless}
\end{axis}
\end{tikzpicture}

\omcaption{The loop of the 2012 question (sum the elements $\ge 128$ of 32\,768 random bytes) with a branch and with a
mask, on sorted and shuffled data. Measured on a laptop (Intel Core Ultra 7 155H) under WSL2, no isolated cores,
pinned to one CPU, best of 30 runs. Data: \texttt{bench\_cpu.py}.}\label{fig:ll:cpu-microarchitecture:branch}
\end{omfigure}

The trading lesson is not ``remove branches''. Most branches on a hot path are predictable (the message type is
usually the same, the risk check usually passes) and cost almost nothing. The expensive ones depend on the market:
which side a message is on, whether a price improves the book. They mispredict in bursts, exactly when the market is
busy, and they are worth measuring one by one.

\section{Frequency scaling, turbo and idle states}

\begin{definition}[Dynamic frequency scaling, idle state]\label{def:ll:cpu-microarchitecture:dvfs}
\emph{Dynamic frequency scaling}\index{dynamic frequency scaling} changes a core's clock (and voltage) while it runs,
according to load, temperature, power budget and the number of active cores; the highest frequencies, called turbo,
are available only while the package stays within its limits. An \emph{idle state}\index{idle state} (a C-state) is a
low-power state a core enters when it has nothing to run, from which it takes microseconds to wake, deeper states
taking longer.
\end{definition}

\begin{definition}[Time-stamp counter]\label{def:ll:cpu-microarchitecture:tsc}
The \emph{time-stamp counter}\index{time-stamp counter} (TSC) is a 64-bit register of each x86 core that counts at a
constant rate on current processors, whatever the core's clock, read by the instructions \texttt{rdtsc} and
\texttt{rdtscp} in a few nanoseconds.
\end{definition}

The two definitions interact. The TSC counts time, not cycles: on this laptop it ticks at about 3.0~GHz while the core
runs at anything from its idle clock to several gigahertz. A benchmark that reports TSC ticks as ``cycles'' is off by
the ratio. The harness reports nanoseconds and, when cycles matter, measures the core's clock with a chain of dependent
register additions, one cycle each. On a server the lesson is operational: a core that has been idle for a while
answers its first message from a deep \omterm{def:ll:cpu-microarchitecture:dvfs}{idle state} and at a low clock, so a trading host limits \omterm{def:ll:cpu-microarchitecture:dvfs}{idle states} and fixes
its frequency (chapter~13).

\begin{dated}{2026-09}{The laptop this book measures on}\label{dat:ll:cpu-microarchitecture:laptop}
The Intel Core Ultra 7 155H, launched in the fourth quarter of 2023, has 16 cores: 6 performance cores with two
hardware threads each, 8 efficient cores and 2 low-power efficient cores, 22 threads in all, with maximum turbo
frequencies of 4.8~GHz on the performance cores and 3.8~GHz on the efficient cores, and 24~MB of shared cache (Intel's
specification page).
\end{dated}

\section{Simultaneous multithreading and heterogeneous cores}

\begin{definition}[Simultaneous multithreading, heterogeneous cores]\label{def:ll:cpu-microarchitecture:smt}
In \emph{simultaneous multithreading}\index{simultaneous multithreading} (SMT, hyper-threading) one core runs two
hardware threads, sharing its execution units, caches and predictors; each appears to the operating system as a CPU.
A processor has \emph{heterogeneous cores}\index{heterogeneous cores} when its cores differ in design, typically large
performance cores and small efficient cores, with different speeds for the same code.
\end{definition}

Both matter to a latency-sensitive thread. Its sibling hardware thread competes for the same execution units and
first-level caches, so trading hosts either disable SMT or leave the sibling of every hot core idle. On a
heterogeneous processor the same code runs at different speeds depending on where the operating system places it; a
hot thread must be pinned to a performance core.

This laptop adds a layer that servers in a colocation site usually do not have: it runs Linux inside a virtual
machine. The 22 CPUs Linux sees are virtual CPUs, which the hypervisor schedules on whatever physical core it chooses,
and the processor's report of its own core types is hidden from the guest. \Cref{fig:ll:cpu-microarchitecture:freq}
measures the clock seen by each virtual CPU: the performance and efficient cores do not appear as two groups, and
pinning a thread to a virtual CPU need not pin it to one physical core. The measurements of this book are therefore the laptop's
best-effort figures, and each caption says so.

\begin{omfigure}
\begin{tikzpicture}
\begin{axis}[width=12cm,height=5cm,ybar,bar width=6pt,ymin=0,ymax=5.5,xmin=-0.8,xmax=21.8,
  xtick={0,2,4,6,8,10,12,14,16,18,20},xlabel={virtual CPU},ylabel={measured clock (GHz)},
  tick label style={font=\small},label style={font=\small},grid=major,grid style={gray!25},table/col sep=comma]
\addplot[fill=omDef!55,draw=omDef] table[x=cpu,y=ghz_reg]{figdata/low-latency/02-cpu-microarchitecture/measured_freq.csv};
\end{axis}
\end{tikzpicture}

\omcaption{The clock each of the laptop's 22 virtual CPUs runs a chain of dependent register additions at (one
addition a cycle), best of 15 runs. The hypervisor hides which physical core runs which virtual CPU; the six
performance and ten efficient cores do not appear as groups. Measured on a laptop (Intel Core Ultra 7 155H) under WSL2,
no isolated cores, the machine otherwise idle. Data: \texttt{bench\_cpu.py}.}\label{fig:ll:cpu-microarchitecture:freq}
\end{omfigure}

\section{Tutorial: timing a kernel honestly}\label{tut:ll:cpu-microarchitecture:tut}

\begin{tutorial}
\textbf{Goal.} Build the \omterm{def:ll:cpu-microarchitecture:ubench}{microbenchmark} harness, then use it to measure a \omterm{def:ll:cpu-microarchitecture:branch}{branch misprediction}, a dependency chain and
the clock. \textbf{End state:}
\cref{fig:ll:cpu-microarchitecture:branch,fig:ll:cpu-microarchitecture:chains,fig:ll:cpu-microarchitecture:freq} and an
estimate of the misprediction penalty in cycles.
\begin{tutsteps}
\item \textbf{Keep the code alive and fence the counter.} An empty \texttt{asm} statement that takes the value as an
input is a \omterm{def:ll:cpu-microarchitecture:ubench}{compiler barrier}: the optimiser must compute the value and cannot move work across it. The timed region
starts with \texttt{lfence; rdtsc} and ends with \texttt{rdtscp; lfence}, so that earlier work finishes before the
first read and the second read waits for the region.
\omcode{code/firm/ubench/cpp/firm_ubench.hpp}{22}{41}{The C++20 harness: barriers and fenced reads of the time-stamp counter.}\label{lst:ll:cpu-microarchitecture:cpp}
\item \textbf{The same in Rust.} Rust has \texttt{std::hint::black\_box} for the barrier; the counter needs the
intrinsics of \texttt{std::arch} and an \texttt{unsafe} block, whose safety argument is a comment.
\omcode{code/firm/ubench/rust/src/lib.rs}{7}{27}{The Rust twin's fenced reads.}\label{lst:ll:cpu-microarchitecture:rust}
\item \textbf{The kernels.} The branchy loop keeps its branch because an opaque \texttt{asm} statement sits inside it;
without it, the compiler may turn the \texttt{if} into a conditional move and the experiment disappears.
\omcode{code/low-latency/02-cpu-microarchitecture/cpp/ll_cpu.hpp}{9}{31}{The 2012 loop with a branch and with a mask.}\label{lst:ll:cpu-microarchitecture:kernels}
\item \textbf{Run} \texttt{python bench\_cpu.py}: it compiles \texttt{ll\_cpu\_bench.cpp}, runs the three modes pinned
with \texttt{taskset}, and writes the measured CSVs and their \texttt{.meta} files; \texttt{ll\_cpu.penalty()} turns
the branch figures into a penalty in cycles.
\end{tutsteps}
\textbf{What to change next.} Remove the \texttt{asm} statement from \texttt{sum\_branchy}, look at the disassembly
(\texttt{objdump -d}) and time it again; replace the random bytes by a pattern of period 8 and watch the predictor
learn it.
\end{tutorial}

\section{Build: the microbenchmark harness}\label{bld:ll:cpu-microarchitecture:ubench}

\begin{build}
\bfield{Purpose} One way to time code across the book, in C++20, Rust and the Python drivers that produce every
measured figure, so that numbers from different chapters are comparable and say where they came from.
\bfield{Interface} C++20 \texttt{firm::ubench}: \texttt{do\_not\_optimize(x)}, \texttt{clobber()},
\texttt{tsc\_start()}, \texttt{rdtscp()}, \texttt{Clock::calibrate(ms)} and \texttt{Clock::ns(ticks)},
\texttt{overhead()}, \texttt{run(fn, reps, warm)}, \texttt{quantile(v, p)}. Rust \texttt{firm\_ubench}: the same
functions. Python \texttt{firm\_ubench}: \texttt{compile\_cpp}, \texttt{run(exe, *args, cpus)},
\texttt{machine()}, \texttt{write\_measured(csv, header, rows, **meta)}, \texttt{quantiles}.
\bfield{Rules} Every timed region is fenced; results are nanoseconds (ticks only inside the harness); every measured CSV
carries a \texttt{.meta} file naming the machine, compiler, flags, CPU list and load; drivers are named
\texttt{bench\_*.py} and are never run by \texttt{make figdata}.
\bfield{Acceptance tests} \texttt{code/firm/ubench/}: calibration between 0.5 and 6 ticks per nanosecond, a monotonic
counter, work costing more than an empty region (C++, Rust), the driver's CSV and metadata round trip (Python).
\bfield{Stretch} Read the core's cycle counter through \texttt{perf\_event\_open} where it is permitted; report the
spread of repeated runs with each figure.
\end{build}

\begin{omsources}
\item Stack Overflow, ``Why is processing a sorted array faster than processing an unsorted array?'', question
11227809, 2012.
\item A.~Fog, \emph{The microarchitecture of Intel, AMD and VIA CPUs}, updated 2026.
\item Chips and Cheese, ``Popping the hood on Golden Cove'' (2021) and ``Lion Cove: Intel's P-core roars'' (2024).
\item Intel, \emph{Intel 64 and IA-32 Architectures Optimization Reference Manual}.
\end{omsources}

\section{Exercises}

\begin{exercise}[$\star$]\label{exo:ll:cpu-microarchitecture:1}
A core retires 2.4 instructions a cycle at 4.0~GHz. How many instructions does a hot path of \qty{600}{\nano\second}
execute?
\end{exercise}

\begin{exercise}[$\star$]\label{exo:ll:cpu-microarchitecture:2}
A benchmark reports 3\,000 TSC ticks for a function on a laptop whose TSC ticks at 3.0~GHz while the core runs at
4.5~GHz. How long did the function take, and how many core cycles?
\end{exercise}

\begin{exercise}[$\star$]\label{exo:ll:cpu-microarchitecture:3}
An operation has a latency of 4 cycles and the core can start two per cycle. How many independent chains does a loop
need to reach the core's throughput, and what fraction of it does a loop with three chains reach?
\end{exercise}

\begin{exercise}[$\star\star$]\label{exo:ll:cpu-microarchitecture:4}
A message handler has a branch on the message's side that mispredicts 30\% of the time, with a penalty of 20 cycles. A
branchless version costs 4 extra cycles every message. Which is faster, and by how much per message?
\end{exercise}

\begin{exercise}[$\star\star$]\label{exo:ll:cpu-microarchitecture:5}
From \cref{fig:ll:cpu-microarchitecture:branch}, estimate the misprediction penalty in nanoseconds, then in cycles at
the clock of \cref{fig:ll:cpu-microarchitecture:freq} for CPU~2. What assumption about the shuffled data does the
estimate rest on?
\end{exercise}

\begin{exercise}[$\star\star$]\label{exo:ll:cpu-microarchitecture:6}
Why can the \omterm{def:ll:cpu-microarchitecture:ooo}{reorder buffer} not hide a dependency chain, however large it is? What does it hide instead?
\end{exercise}

\begin{exercise}[$\star\star\star$]\label{exo:ll:cpu-microarchitecture:7}
\emph{Coding.} Add a mode to \texttt{ll\_cpu\_bench.cpp} that runs the branchy loop on data whose values alternate in
blocks of $b$ elements above and below 128, for $b = 1, 2, 4, 8, 64$. What does the predictor learn, and at which $b$
does the loop reach the sorted speed?
\end{exercise}

\begin{exercise}[$\star\star\star$]\label{exo:ll:cpu-microarchitecture:8}
\emph{Find the flaw.} ``We measured our order encoder at 40 cycles with a loop of \texttt{add \$1} instructions as a
cycle reference, on a warm core, calling it ten million times in a row on the same order.''
\end{exercise}

\section{Problem: The Branch That Cost a Quote}

\begin{problem}[{Weekend problem --- one \texttt{if} in the quoting path}]\label{pb:ll:cpu-microarchitecture:1}
A quoting engine's hot path has one data-dependent branch: whether an incoming trade was on the bid or the ask. On a
quiet day the sides alternate in long runs; on a busy day they are close to random. The team measures the 2012 loop as
a stand-in (\cref{fig:ll:cpu-microarchitecture:branch}) and must decide whether to rewrite the branch.

\medskip\noindent\textbf{Part I --- The measurement.}
\begin{enumerate}
\item What are the four measured costs per element?
\item Why is the sorted branchy loop faster than the branchless one?
\item Estimate the misprediction penalty in nanoseconds.
\item Convert it to cycles with the measured clock of CPU~2.
\end{enumerate}

\medskip\noindent\textbf{Part II --- The model.}
\begin{enumerate}[resume]
\item Write the expected cost per iteration of the branchy and branchless loops as functions of the misprediction
rate.
\item At what misprediction rate do they cost the same?
\item On a quiet day the branch mispredicts on 1\% of trades. Which version wins, by how much per trade?
\item On a busy day it mispredicts on 35\%. Which wins, by how much?
\end{enumerate}

\medskip\noindent\textbf{Part III --- What it is worth.}
\begin{enumerate}[resume]
\item On the busy day trades arrive at 50\,000 a second in bursts. How much processing time a second does the branchy
version waste?
\item Compared with the path's budget of \qty{1.3}{\micro\second} for software stages (chapter~1), is the penalty per
mispredicted trade large or small?
\item Why do mispredictions cluster on the days that matter?
\item What would a predictor that learnt the busy day's pattern change?
\end{enumerate}

\medskip\noindent\textbf{Part IV --- The decision.}
\begin{enumerate}[resume]
\item State the \emph{named result}: the misprediction penalty measured on this laptop, in nanoseconds and cycles, and the
misprediction rate above which the branchless rewrite wins.
\item Why is the measured penalty above the published figure for the core family?
\item What does the TSC measure in this experiment, and why is it not cycles?
\item Why pin the benchmark to one CPU, and what does pinning mean under a hypervisor?
\item What would a compiler do to the branchy loop if the \texttt{asm} statement were removed?
\item Name another branch on a trading hot path that is predictable, and one that is not.
\item How would you measure the misprediction rate on the production host?
\item In one sentence: when is branchless code worth writing?
\end{enumerate}
\end{problem}

\section{Interview questions}

\begin{interviewq}[$\star$ \iqroles{developer}]\label{iq:ll:cpu-microarchitecture:1}
Why can sorting an array make a loop over it several times faster, when the loop's work does not depend on the order?
\end{interviewq}

\begin{interviewq}[$\star\star$ \iqroles{developer}]\label{iq:ll:cpu-microarchitecture:2}
Explain \omterm{def:ll:cpu-microarchitecture:ooo}{out-of-order execution}. What limits how much independent work the core can find?
\end{interviewq}

\begin{interviewq}[$\star\star$ \iqroles{developer}]\label{iq:ll:cpu-microarchitecture:3}
Why does summing an array with four accumulators run faster than with one? Why does the compiler not do it for you?
\end{interviewq}

\begin{interviewq}[$\star\star$ \iqroles{developer}]\label{iq:ll:cpu-microarchitecture:4}
You want to time a function that takes 50 nanoseconds. How do you do it, and what can go wrong?
\end{interviewq}

\begin{interviewq}[$\star\star$ \iqroles{developer}]\label{iq:ll:cpu-microarchitecture:5}
Should a latency-critical thread share a core with its hyper-threading sibling? Why?
\end{interviewq}

\begin{interviewq}[$\star\star\star$ \iqroles{developer}]\label{iq:ll:cpu-microarchitecture:6}
The first order after a quiet minute is always slower than the next ones. Name three causes in the processor and how
you would confirm each.
\end{interviewq}
